\documentclass[a4paper,11pt]{article}

% =========================================================================
%                                PAQUETES
% =========================================================================

% --- IDIOMA Y CODIFICACIÓN ---
\usepackage[utf8]{inputenc}
\usepackage[T1]{fontenc}
\usepackage[spanish,es-tabla,es-nodecimaldot]{babel}

% --- MATEMÁTICAS Y FÍSICA ---
\usepackage{amsmath, amsfonts, amssymb, mathtools}
\usepackage{amsthm} 
\usepackage{cancel} 
\usepackage{esint}  
\usepackage{empheq} 
\usepackage{physics}
\usepackage{siunitx}
\usepackage{subcaption}
\usepackage{caption}

% --- DISEÑO Y GRÁFICOS ---
\usepackage{geometry}
\geometry{left=2.5cm, right=2.5cm, top=3cm, bottom=3cm}
\usepackage{graphicx}
\usepackage{float}    
\usepackage[table]{xcolor}
\usepackage{tikz}     
\usetikzlibrary{babel, arrows.meta, calc, positioning, 3d, decorations.markings, patterns, shadings, angles, quotes, matrix}
\usepackage{pgfplots}
\pgfplotsset{compat=1.17}
\usepackage[most]{tcolorbox}
\tcbuselibrary{skins, breakable}
\usepackage{array}
\usepackage{textpos}
\usepackage{setspace}
\usepackage{multirow}

% --- CONFIGURACIÓN DE PÁRRAFO, TABLAS E INTERLINEADO ---
\usepackage{cellspace}
\setlength{\parskip}{0.6em}
\setlength{\parindent}{0pt}
\setstretch{1.15} % Interlineado ampliado para un aspecto más limpio y formal
\renewcommand{\arraystretch}{1.5} 
\setlength\cellspacetoplimit{10pt}   
\setlength\cellspacebottomlimit{10pt}

% --- ENCABEZADO Y PIE DE PÁGINA ---
\usepackage{fancyhdr}
\pagestyle{fancy}
\fancyhf{} 
\fancyhead[L]{\textit{Laboratorio de Mecánica y Ondas}}
\fancyhead[R]{Práctica 8: Efecto Doppler}
\fancyfoot[C]{\thepage}
\renewcommand{\headrulewidth}{0.4pt}

% --- SECCIONES Y SUBSECCIONES ---
\setcounter{secnumdepth}{2}
\setcounter{tocdepth}{2}

% --- NAVEGACIÓN ---
\usepackage{hyperref}
\hypersetup{
    colorlinks=true,
    linkcolor=blue!70!black,
    citecolor=blue,
    urlcolor=blue
}

% =========================================================================
%                        ENTORNOS PERSONALIZADOS
% =========================================================================
% Estética minimalista de artículo científico: líneas laterales, tonos sobrios y sin bordes cerrados.

\newtcolorbox{resultado}[1]{
    enhanced,
    breakable,
    frame hidden,
    interior hidden,
    boxrule=0pt,
    leftrule=3pt,
    arc=0pt,
    colframe=blue!60!black,
    colback=blue!2!white,
    fonttitle=\bfseries\sffamily,
    coltitle=blue!60!black,
    title={#1},
    before skip=14pt,
    after skip=14pt,
    left=12pt, right=10pt, top=4pt, bottom=4pt,
    attach title to upper=\par\vspace{0.4em}
}

\newtcolorbox{nota}[1]{
    enhanced,
    breakable,
    frame hidden,
    interior hidden,
    boxrule=0pt,
    leftrule=2pt,
    arc=0pt,
    colframe=gray!70!blue,
    colback=white,
    fonttitle=\bfseries\sffamily,
    coltitle=gray!70!blue,
    title={Nota: #1},
    before skip=12pt,
    after skip=12pt,
    left=12pt, right=10pt, top=4pt, bottom=4pt,
    attach title to upper=\par\vspace{0.3em}
}

\newtcolorbox{referencia}[1]{
    enhanced,
    breakable,
    frame hidden,
    interior hidden,
    boxrule=0pt,
    leftrule=3pt,
    arc=0pt,
    colframe=blue!70!black,
    colback=blue!3!white,
    fonttitle=\bfseries\sffamily,
    coltitle=blue!70!black,
    fontupper=\itshape, % Fuerza la cursiva en todo el texto del entorno
    title={#1},
    before skip=14pt,
    after skip=14pt,
    left=12pt, right=12pt, top=6pt, bottom=6pt,
    attach title to upper=\par\vspace{0.3em}
}

% =========================================================================
%                     VARIABLES PARA LA PORTADA
% =========================================================================

\newcommand{\Universidad}{Universidad de Salamanca}
\newcommand{\LogoUniversidad}{logos-universidad/USAL}
\newcommand{\Facultad}{Facultad de Ciencias}
\newcommand{\LogoFacultad}{logos-universidad/logo}
\newcommand{\Titulacion}{Grado en Física}
\newcommand{\titulo}{Laboratorio de Mecánica y Ondas\\Práctica 8: Efecto Doppler}
\newcommand{\tipotrabajo}{Informe de laboratorio}
\newcommand{\miNombre}{David Gómez García \\ Daniel Martín Guerrero}
\newcommand{\miTutor}{Mario Amado \\ Ivan de Martino \\ Manuel Pino}
\newcommand{\miUbicacion}{Salamanca}

% =========================================================================
%                                DOCUMENTO
% =========================================================================

\setlength{\headheight}{13.6pt}
\begin{document}

% --- PORTADA ---
%%%%%%%%%%%%%%%%%%%%%%%%%%%%%%%%%%%%%%%%%%%%%%%%%%%%%%%%%%%%%%%%%%%%%%%%
% Portada — estilo artículo científico, motivo: frentes de onda Doppler
% Requiere (definidos en main.tex): colores azulprof, azulacero, azulclaro,
% grispizarra, grisclaro; y las macros \Universidad, \Facultad, \Titulacion,
% \Asignatura, \NumPractica, \TituloPractica, \tipotrabajo, \miNombre,
% \miTutor, \miUbicacion.
%%%%%%%%%%%%%%%%%%%%%%%%%%%%%%%%%%%%%%%%%%%%%%%%%%%%%%%%%%%%%%%%%%%%%%%%

\begin{titlepage}
\thispagestyle{empty}
\pagecolor{white}

\begin{tikzpicture}[remember picture, overlay,
    pg/.style={x=1cm, y=1cm}]
  \coordinate (O) at (current page.south west);

  % ----------------------------------------------------------------
  % Motivo gráfico: frentes de onda de una fuente que se desplaza hacia
  % la derecha (compresión a la derecha, dilatación a la izquierda).
  % ----------------------------------------------------------------
  \begin{scope}
    \clip (O) rectangle ($(O)+(21cm,24.6cm)$);
    \foreach \k in {0,...,15}{
      \pgfmathsetmacro{\cx}{8.2+0.36*\k}
      \pgfmathsetmacro{\rr}{8.4-0.60*\k}
      \pgfmathsetmacro{\pc}{14+5.2*\k}
      \pgfmathsetmacro{\lw}{0.25+0.035*\k}
      \draw[azulprof!\pc!azulclaro, line width=\lw pt]
        ($(O)+(\cx cm,17.2cm)$) circle (\rr cm);
    }
    \fill[azulprof] ($(O)+(8.2cm+0.36cm*15,17.2cm)$) circle (2.2pt);
  \end{scope}

  % ----------------------------------------------------------------
  % Señal con desplazamiento Doppler (chirp): frecuencia creciente
  % ----------------------------------------------------------------
  \begin{scope}[shift={(O)}]
    \draw[azulacero, line width=0.9pt, domain=2.2:18.8, samples=420, smooth]
      plot ({\x}, {12.95 + 0.33*sin(deg((\x-2.2)*(2.2+0.58*(\x-2.2))))});
  \end{scope}

  % ----------------------------------------------------------------
  % Cabecera: logotipos y filete
  % ----------------------------------------------------------------
  \node[anchor=north west, inner sep=0] at ($(O)+(2.2cm,27.9cm)$)
    {\includegraphics[height=1.42cm]{Logo USAL.png}};
  \node[anchor=north east, inner sep=0] at ($(O)+(18.8cm,27.9cm)$)
    {\includegraphics[height=1.42cm]{Logo Facultad de Ciencias.png}};
  \draw[grispizarra!60, line width=0.4pt]
    ($(O)+(2.2cm,26.0cm)$) -- ($(O)+(18.8cm,26.0cm)$);
  \node[anchor=north west, inner sep=0, text=grispizarra, font=\scshape\small]
    at ($(O)+(2.2cm,25.75cm)$)
    {\textls[80]{\Universidad\ · \Facultad}};
  \node[anchor=north east, inner sep=0, text=grispizarra, font=\scshape\small]
    at ($(O)+(18.8cm,25.75cm)$)
    {\textls[80]{\Titulacion}};

  % ----------------------------------------------------------------
  % Bloque del título
  % ----------------------------------------------------------------
  \node[anchor=south west, inner sep=0, text=azulacero, font=\scshape\large,
        text width=16.6cm]
    at ($(O)+(2.2cm,11.2cm)$)
    {\textls[120]{\Asignatura\ ·\ \NumPractica}};

  \node[anchor=north west, inner sep=0, text=azulprof, text width=16.6cm,
        align=left]
    at ($(O)+(2.2cm,11.0cm)$)
    {\fontsize{35}{39}\selectfont\bfseries\TituloPractica};

  \draw[azulacero, line width=1.4pt]
    ($(O)+(2.2cm,7.35cm)$) -- ($(O)+(4.7cm,7.35cm)$);

  \node[anchor=north west, inner sep=0, text=grispizarra, text width=16.6cm,
        align=left, font=\itshape\large]
    at ($(O)+(2.2cm,7.05cm)$)
    {$\displaystyle f_{\mathrm{R}} \;=\; f_{\mathrm{E}}\,
      \frac{v_{\mathrm{s}}-v_{\mathrm{R}}}{v_{\mathrm{s}}-v_{\mathrm{E}}}$};

  % ----------------------------------------------------------------
  % Pie: tipo de trabajo, autores, tutores, lugar y fecha
  % ----------------------------------------------------------------
  \draw[grispizarra!60, line width=0.4pt]
    ($(O)+(2.2cm,5.1cm)$) -- ($(O)+(18.8cm,5.1cm)$);

  \node[anchor=north west, inner sep=0, text width=4.6cm, align=left]
    at ($(O)+(2.2cm,4.75cm)$)
    {{\scshape\small\color{azulacero}\textls[80]{Documento}}\\[3pt]
     {\large\bfseries\color{azulprof}\tipotrabajo}};

  \node[anchor=north west, inner sep=0, text width=5.6cm, align=left]
    at ($(O)+(7.2cm,4.75cm)$)
    {{\scshape\small\color{azulacero}\textls[80]{Autores}}\\[3pt]
     {\large\color{black!85}\miNombre}};

  \node[anchor=north west, inner sep=0, text width=5cm, align=left]
    at ($(O)+(13.8cm,4.75cm)$)
    {{\scshape\small\color{azulacero}\textls[80]{Tutores}}\\[3pt]
     {\large\color{black!85}\miTutor}};

  \node[anchor=south east, inner sep=0, text=grispizarra, font=\scshape\small]
    at ($(O)+(18.8cm,1.45cm)$)
    {\textls[100]{\miUbicacion, \number\year}};
  \fill[azulprof] ($(O)+(0,0)$) rectangle ($(O)+(21cm,0.45cm)$);
  \fill[azulacero] ($(O)+(0,0)$) rectangle ($(O)+(7cm,0.45cm)$);
  \fill[azulclaro] ($(O)+(0,0)$) rectangle ($(O)+(3cm,0.45cm)$);
\end{tikzpicture}

\end{titlepage}
 
\restoregeometry

% --- ÍNDICE ---
\newpage
\tableofcontents
\newpage

% --- CONTENIDO ---
\section{Introducción}
La experiencia y su posterior análisis recogidos en este informe presentan como temática nuclear un fenómeno físico conocido como efecto Doppler.

\subsection{Objetivos}
\begin{itemize}
    \item Hallar la velocidad con la que se propagan ondas ultrasonidos en un medio como el aire.
    \item Evidenciar y calcular las diferencias entre la frecuencia emitida y recibida ante varias situaciones dinámicas del emisor y receptor.
\end{itemize}

\subsection{Material empleado}
Para el correcto desarrollo del montaje experimental, se ha hecho uso del siguiente instrumental:

\vspace{0.4cm} 
\begin{figure}[H]
    % --- Columna izquierda para la lista y la explicación ---
    \begin{minipage}[t]{0.35\textwidth}
        \vspace{0pt}
        \begin{itemize}
            \item Unidad de emisión y detección de ultrasonidos.
            \item Carril y coche autopropulsado.
            \item Emisor, receptor y soportes.
            \item Barrera fotoeléctrica y temporizador.
            \item Ordenador y unidad de medida Cobra-3.
        \end{itemize}
        
        \vspace{0.4cm}
        
        
    \end{minipage}%
    \hfill 
    % --- Columna derecha para la imagen con pie de foto reducido ---
    \begin{minipage}[t]{0.6\textwidth}
        \vspace{0pt}
        \centering
        \includegraphics[width=\textwidth]{} 
        \caption{Montaje experimental empleado.}
        \label{fig:montaje_experimental}
    \end{minipage}
\end{figure}

\subsection{Fundamentos teóricos}

\begin{referencia}{Efecto Doppler}
El efecto Doppler es un fenómeno físico que describe el cambio en la frecuencia percibida de una onda cuando la fuente emisora de la onda y el observador están en movimiento relativo 
entre sí. Este cambio de frecuencia no se debe a una modificación real de la frecuencia de la 
onda emitida, sino a un cambio en la forma en que el observador percibe la onda debido al 
movimiento. 
La causa del efecto Doppler radica en la compresión o dilatación de las ondas debida al 
movimiento relativo entre la fuente y el observador. Si la fuente se mueve hacia el 
observador, las ondas se comprimen, lo que resulta en una longitud de onda más corta y una 
frecuencia percibida más alta. Si la fuente se aleja, las ondas se dilatan, lo que resulta en una 
longitud de onda más larga y una frecuencia percibida más baja. El efecto Doppler se aplica 
a cualquier tipo de onda, incluyendo ondas sonoras, ondas de luz, ondas de agua, etc. 

Si la fuente se acerca al observador ($v_s>0$), la longitud de onda se acorta y la frecuencia se hace mayor. En ondas de sonido se aprecia una frecuencia más aguda. Por el contrario, si el emisor se aleja del receptor ($v_s<0$), la longitud de onda 
aparente es mayor y la frecuencia menor. El sonido se hace más grave. 
\cite{FV6-ECJ}
\begin{figure}[H]
    \centering
    \resizebox{\textwidth}{!}{%
    \begin{tikzpicture}[font=\sffamily, >=stealth, scale=0.85, every node/.style={scale=0.9}]
        
        % ==========================================
        % PANEL IZQUIERDO: FRENTES DE ONDA ESPACIALES
        % ==========================================
        \begin{scope}[shift={(0,0)}]
            % Eje espacial
            \draw[->, thick, gray!80] (-5.5, 0) -- (4.5, 0) node[right, text=black] {$x$};
            
            % Definición de radios y centros (simulando v_s y v_onda)
            \def\rA{3.8} \def\cA{0}
            \def\rB{2.8} \def\cB{0.7}
            \def\rC{1.8} \def\cC{1.4}
            \def\rD{0.8} \def\cD{2.1}
            
            % Trazado de los frentes de onda (circunferencias)
            \draw[thick, gray!70] (\cA,0) circle (\rA);
            \draw[thick, gray!70] (\cB,0) circle (\rB);
            \draw[thick, gray!70] (\cC,0) circle (\rC);
            \draw[thick, gray!70] (\cD,0) circle (\rD);
            
            % Posiciones sucesivas de la fuente
            \fill[teal!80!black] (\cA,0) circle (2.5pt) node[below=2pt] {$x_0$};
            \fill[teal!80!black] (\cB,0) circle (2.5pt) node[below=2pt] {$x_1$};
            \fill[teal!80!black] (\cC,0) circle (2.5pt) node[below=2pt] {$x_2$};
            \fill[teal!80!black] (\cD,0) circle (2.5pt) node[below=2pt] {$x_3$};
            \node[teal!80!black, above=4pt] at (\cA,0) {\textbf{S}};
            
            % Receptores
            \fill[magenta!80!black] (-4.5,0) circle (3pt) node[below=4pt] {$R_1$};
            \fill[blue!80!black] (3.5,0) circle (3pt) node[below=4pt] {$R_2$};
            
            % Vectores de radio
            \draw[->, dashed, gray!90] (\cA,0) -- ++(110:\rA) node[midway, left=2pt, fill=white, inner sep=1pt, text=black!80] {$r_0$};
            \draw[->, dashed, gray!90] (\cB,0) -- ++(80:\rB) node[midway, right=1pt, fill=white, inner sep=1pt, text=black!80] {$r_1$};
            \draw[->, dashed, gray!90] (\cC,0) -- ++(50:\rC) node[midway, right=1pt, fill=white, inner sep=1pt, text=black!80] {$r_2$};
            \draw[->, dashed, gray!90] (\cD,0) -- ++(20:\rD) node[midway, right=1pt, fill=white, inner sep=1pt, text=black!80] {$r_3$};
            
            % Ecuaciones paramétricas (Radios)
            \node[anchor=east, align=right, text=black!80] at (-5.5, 2.5) {
                $r_0 = v(t - t_0)$\\
                $r_1 = v(t - t_1)$\\
                $r_2 = v(t - t_2)$\\
                $r_3 = v(t - t_3)$
            };
            
            % Ecuaciones paramétricas (Posiciones de la fuente)
            \node[anchor=east, align=right, text=black!80] at (-5.5, -2.2) {
                $x_0 = v_s t_0$\\
                $x_1 = v_s t_1$\\
                $x_2 = v_s t_2$\\
                $x_3 = v_s t_3$
            };
        \end{scope}

        % ==========================================
        % PANEL DERECHO: GRÁFICAS TEMPORALES
        % ==========================================
        \begin{scope}[shift={(6.5, 0)}]
            \def\w{5.5}  % Ancho de las gráficas
            \def\ha{1.2} % Semiamplitud
            
            % Gráfica 1: Onda emitida por S
            \begin{scope}[shift={(0, 3)}]
                \draw[black!15, thin, step=0.5] (0, -\ha) grid (\w, \ha);
                \draw[->, thick, gray!80] (0,0) -- (\w+0.5, 0);
                \draw[->, thick, gray!80] (0,-\ha-0.2) -- (0, \ha+0.2) node[left, text=black!80] {Amp.};
                \draw[teal!80!black, thick, domain=0:\w, samples=150] plot (\x, {0.9*sin(\x*180*1.5)});
                \node[draw=teal!80!black, fill=white, text=teal!80!black, rounded corners=2pt, inner sep=3pt] at (\w-0.5, \ha-0.2) {\textbf{S}};
            \end{scope}
            
            % Gráfica 2: Onda percibida por R1 (Frecuencia menor)
            \begin{scope}[shift={(0, 0)}]
                \draw[black!15, thin, step=0.5] (0, -\ha) grid (\w, \ha);
                \draw[->, thick, gray!80] (0,0) -- (\w+0.5, 0);
                \draw[->, thick, gray!80] (0,-\ha-0.2) -- (0, \ha+0.2) node[left, text=black!80] {Amp.};
                \def\start{1.8} % Retardo de llegada
                \draw[magenta!80!black, thick] (0,0) -- (\start, 0);
                \draw[magenta!80!black, thick, domain=\start:\w, samples=100] plot (\x, {0.9*sin((\x-\start)*180*0.9)});
                \node[draw=magenta!80!black, fill=white, text=magenta!80!black, rounded corners=2pt, inner sep=3pt] at (\w-0.5, \ha-0.2) {\textbf{$R_1$}};
            \end{scope}
            
            % Gráfica 3: Onda percibida por R2 (Frecuencia mayor)
            \begin{scope}[shift={(0, -3)}]
                \draw[black!15, thin, step=0.5] (0, -\ha) grid (\w, \ha);
                \draw[->, thick, gray!80] (0,0) -- (\w+0.5, 0) node[below, text=black] {Tiempo};
                \draw[->, thick, gray!80] (0,-\ha-0.2) -- (0, \ha+0.2) node[left, text=black!80] {Amp.};
                \def\start{0.8} % Retardo de llegada
                \draw[blue!80!black, thick] (0,0) -- (\start, 0);
                \draw[blue!80!black, thick, domain=\start:\w, samples=150] plot (\x, {0.9*sin((\x-\start)*180*2.5)});
                \node[draw=blue!80!black, fill=white, text=blue!80!black, rounded corners=2pt, inner sep=3pt] at (\w-0.5, \ha-0.2) {\textbf{$R_2$}};
            \end{scope}
        \end{scope}
    \end{tikzpicture}
    }%
    \caption{Esquema del efecto Doppler generado por una fuente \textbf{S} en movimiento a velocidad $v_s$. A la izquierda se representan los frentes de onda espaciales emitidos en distintos instantes, lo que evidencia la compresión en la dirección del movimiento y la dilatación en la opuesta. A la derecha, las gráficas temporales comparan la onda emitida con las señales desfasadas percibidas por los receptores $R_1$ y $R_2$.}
    \label{fig:doppler_frentes}
\end{figure}
\end{referencia}

\begin{figure}[H]
    \centering
    \begin{tikzpicture}[font=\sffamily, >=stealth, scale=1.2]
        
        % ==========================================
        % EJE DE REFERENCIA ESPACIAL
        % ==========================================
        \draw[->, gray!40, thin] (-2, 0) -- (10, 0) node[right, text=gray!80, scale=0.9] {$x$};
        
        % ==========================================
        % ONDA PRINCIPAL (Nuevo tono: cyan!70!black)
        % ==========================================
        \draw[cyan!70!black, thick, densely dashed, domain=-1.5:1, samples=50] plot (\x, {sin(\x*180)});
        \draw[cyan!70!black, thick, densely dashed, domain=7:9.5, samples=50] plot (\x, {sin(\x*180)});
        
        % Zona continua central
        \draw[cyan!70!black, very thick, domain=1:7, samples=150] plot (\x, {sin(\x*180)});
        
        % ==========================================
        % MARCADORES DE EMISOR Y RECEPTOR
        % ==========================================
        \draw[thick, gray!60, densely dotted] (1, -1.4) -- (1, 1.4);
        \node[draw=gray!50, fill=white, rounded corners=3pt, font=\bfseries\sffamily\small, text=gray!80!black, inner sep=4pt] at (1, 1.7) {emisor};
        
        \draw[thick, gray!60, densely dotted] (7, -1.4) -- (7, 1.4);
        \node[draw=gray!50, fill=white, rounded corners=3pt, font=\bfseries\sffamily\small, text=gray!80!black, inner sep=4pt] at (7, 1.7) {receptor};
        
        % ==========================================
        % VECTORES DE VELOCIDAD RELATIVA
        % ==========================================
        \draw[->, thick, gray!80!black] (1.0, -1.7) -- (0.0, -1.7) node[midway, below] {$v_E - v_s$};
        \draw[->, thick, gray!80!black] (7.0, -1.7) -- (6.0, -1.7) node[midway, below] {$v_R - v_s$};
        
        % ==========================================
        % ACOTACIÓN DE LA LONGITUD DE ONDA
        % ==========================================
        \draw[gray!50, thin] (2.5, 1.1) -- (2.5, 1.6);
        \draw[gray!50, thin] (4.5, 1.1) -- (4.5, 1.6);
        \draw[|<->|, thick, cyan!70!black] (2.5, 1.45) -- (4.5, 1.45) node[midway, fill=white, inner sep=2pt] {$\lambda$};
        
    \end{tikzpicture}
    \caption{Movimientos relativos del emisor y del receptor en el sistema de referencia de la onda, asumiendo $u > v_E, v_R$. Se toma el sentido positivo de propagación de izquierda a derecha.}
    \label{fig:doppler_relativo}
\end{figure}


\newpage
\section{Procedimiento experimental}
 En primer lugar se ha comenzado la experiencia midiendo el valor de la temperatura del aire, $T= \text{º }C$, a fin de determinar el valor teórico o valor esperado de la velocidad del sonido en el aire ($\tilde{v_s})$.
\begin{equation}
    \tilde{v_{s}} = 331,3 \sqrt{1+\frac{T}{273,15}}\unit{\meter\per\second})
    \label{eq:vs}
\end{equation}
\subsection{Tratamiento de incertidumbres}
\begin{resultado}{Incertidumbre de la lectura térmica}
    

Obtenida la lectura de temperatura, $T$, con una resolución del termómetro $\Delta = 0,1  \text{º }C$, se procede a calcular la incertidumbre asociada a dicha medida:
\[\sigma^2 = \frac{\Delta^2}{12} \xrightarrow{} u(x) = \frac{\Delta}{\sqrt{12}} \approx
0,03  \text{º }C\] 
En consecuencia, el registro térmico al comienzo de la experiencia ha sido $T = t \pm 0,03  \textbf{º }C$.
\end{resultado}
\begin{resultado}{Incertidumbre de la velocidad del sonido}

Determinar la incertidumbre asociada a la lectura de la temperatura es un paso clave antes de poder hallar por propagación la incertidumbre propia de $\tilde{v_{s}}$.
\begin{equation*}
    u^2(\tilde{v_s}) = \left (\frac{\partial \tilde{v_s}}{\partial T}\right )^{\!2} u^2(T)
\end{equation*}
Para calcular la derivada parcial nos apoyamos en la dependencia térmica de la velocidad del sonido en el aire [\ref{eq:vs}].
Realizando la derivada parcial respecto de la temperatura evaluada para el valor experimental:
\begin{equation*}
    \frac{\partial v_s}{\partial T} = \frac{331.3}{2 \cdot 273.15 \sqrt{1 + \frac{T}{273.15}}} = \frac{331.3}{546.3 \sqrt{1 + \frac{T}{273.15}}} 
\end{equation*}

Si se sustituye este resultado en la expresión que permite calcular la incertidumbre, y calculamos su raíz cuadrda:
\begin{equation*}
    u (v_s) = \sqrt{\left (\frac{\partial v_s}{\partial T}\right)^{\!2} u^2(T)} = \sqrt{\left (\frac{\partial v_s}{\partial T}\right)^{\!2} \left(\frac{0,1}{\sqrt{12}}\right)^{\!2}} \approx error \, \unit{\meter\per\second}
\end{equation*}

Luego, se espera que la velocidad del sonido en el aire para las condiciones en las que se ha encontrado el laboratoria sea:

\[\tilde{v_s} = \tilde{vs} \pm error\,\unit{\meter\per\second}\]

\end{resultado}

% --- BIBLIOGRAFÍA ---
\newpage
\bibliographystyle{plain}
\bibliography{referencias}

\end{document}
